\section{Results}
\label{sec:results}

All numbers in this section are computed by scripts directly from the run logs.
Numbers in the text refer to the \splitname\ split; Table~\ref{tab:main} lists every split that has been evaluated.

\subsection{Main comparison (E1)}
\begin{table*}[t]
\centering\small
\resizebox{\textwidth}{!}{\begin{tabular}{lllllllll}
\toprule
split & method & seeds & EM mean$\pm$std & EM 95\% CI & F1 & strict validity & hop eff & mean hops \\
\midrule
dev & M0 oracle & 1 & 0.603$\pm$0.000 & [0.573, 0.634] & 0.669$\pm$0.000 & 1.000$\pm$0.000 & 0.610$\pm$0.000 & 2.05$\pm$0.00 \\
dev & B1 SFT & 1 & 0.263$\pm$0.000 & [0.236, 0.291] & 0.322$\pm$0.000 & n/a & n/a & n/a \\
dev & B2a GRPO & 1 & 0.101$\pm$0.000 & [0.083, 0.120] & 0.147$\pm$0.000 & 0.074$\pm$0.000 & 0.052$\pm$0.000 & 4.00$\pm$0.00 \\
dev & B2b GRPO+path & 1 & 0.074$\pm$0.000 & [0.058, 0.091] & 0.114$\pm$0.000 & 0.085$\pm$0.000 & 0.030$\pm$0.000 & 4.00$\pm$0.00 \\
dev & B3 soft & 1 & 0.139$\pm$0.000 & [0.118, 0.161] & 0.189$\pm$0.000 & 0.076$\pm$0.000 & 0.103$\pm$0.000 & 3.00$\pm$0.00 \\
dev & H1 Hop-PPO & 1 & 0.114$\pm$0.000 & [0.095, 0.134] & 0.155$\pm$0.000 & 0.172$\pm$0.000 & 0.099$\pm$0.000 & 2.04$\pm$0.00 \\
\bottomrule
\end{tabular}
}
\caption{In-domain comparison on HybridQA: exact match (EM) with seed mean $\pm$ std and a 95\% bootstrap interval over questions, token F1, strict path validity, hop efficiency and mean hops (definitions in Section~\ref{sec:setup}). Path metrics use questions with a gold path; B1 has no hops. All path-based methods share the encoder, candidates and frozen answer head, and all trained methods use the same rollout and tuning budgets.}
\label{tab:main}
\end{table*}

\paragraph{Accuracy.}
Full-table SFT (B1) is the most accurate method, with \R{EmBOneSft} EM \CI{BOneSft}, and \hop\ (H1) is \D{BOneSft} EM points behind it (Table~\ref{tab:paired}); we therefore do not find that hop-level selection beats SFT in this regime.
All path-based methods are far below the oracle M0 (\R{EmMZero} EM), which reads the gold path with the same frozen answer head, so evidence selection, not the reader, is the main bottleneck.
\hop\ reaches \R{EmHOneHoppo} EM \CI{HOneHoppo}.
Against answer-only GRPO (B2a, \R{EmBTwoAGrpo}) the difference is \D{BTwoAGrpo} points, an interval that includes zero.
Against GRPO that receives the same gold-path signal at the end of the trace (B2b, \R{EmBTwoBGrpoPath}) the difference is \D{BTwoBGrpoPath}, so distributing path credit over hops helped relative to giving it once at the end (H1b).
The one-shot soft selector B3 (\R{EmBThreeSoft}) is not reliably different from \hop\ (\D{BThreeSoft}).

\paragraph{Paths.}
\hop\ has the highest strict path validity of the learned path methods (\R{ValidHOneHoppo}\%, versus \R{ValidBTwoAGrpo}\% for B2a, \R{ValidBTwoBGrpoPath}\% for B2b and \R{ValidBThreeSoft}\% for B3) and uses \R{HopsHOneHoppo} hops on average, close to the mean gold-path length (\PathLenTrain), while both GRPO baselines run to the 4-hop limit.
Validity is nevertheless low in absolute terms, so H2 (auditable, valid paths) is only partially supported: every answer comes with its evidence by construction, but most paths leave the gold path after the first hop (Figure~\ref{fig:curves}, right).

\begin{table}[t]
\centering\small
\resizebox{\columnwidth}{!}{\begin{tabular}{llllll}
\toprule
split & comparison & n questions & diff EM & 95\% CI & p \\
\midrule
dev & H1 Hop-PPO - M0 oracle & 1000 & -0.501 & [-0.532, -0.471] & 0.0000 \\
dev & H1 Hop-PPO - B1 SFT & 1000 & -0.190 & [-0.219, -0.161] & 0.0000 \\
dev & H1 Hop-PPO - B2a GRPO & 1000 & -0.029 & [-0.041, -0.017] & 0.0000 \\
dev & H1 Hop-PPO - B2b GRPO+path & 1000 & +0.030 & [+0.015, +0.046] & 0.0002 \\
dev & H1 Hop-PPO - B3 soft & 1000 & -0.029 & [-0.050, -0.008] & 0.0046 \\
dev & H1 Hop-PPO - A1 terminal-only & 1000 & -0.032 & [-0.045, -0.020] & 0.0000 \\
dev & H1 Hop-PPO - A2 no-redund. & 1000 & -0.054 & [-0.068, -0.039] & 0.0000 \\
dev & H1 Hop-PPO - A3 no-schema & 1000 & -0.004 & [-0.015, +0.007] & 0.4714 \\
dev & H1 Hop-PPO - A4 unrestricted & 1000 & -0.001 & [-0.013, +0.012] & 0.9298 \\
test & H1 Hop-PPO - M0 oracle & 2045 & -0.476 & [-0.497, -0.455] & 0.0000 \\
test & H1 Hop-PPO - B1 SFT & 2045 & -0.175 & [-0.194, -0.156] & 0.0000 \\
test & H1 Hop-PPO - B2a GRPO & 2045 & -0.020 & [-0.028, -0.012] & 0.0000 \\
test & H1 Hop-PPO - B2b GRPO+path & 2045 & +0.026 & [+0.016, +0.036] & 0.0000 \\
test & H1 Hop-PPO - B3 soft & 2045 & -0.018 & [-0.031, -0.005] & 0.0056 \\
test & H1 Hop-PPO - A1 terminal-only & 2045 & -0.023 & [-0.032, -0.015] & 0.0000 \\
test & H1 Hop-PPO - A2 no-redund. & 2045 & -0.052 & [-0.062, -0.042] & 0.0000 \\
test & H1 Hop-PPO - A3 no-schema & 2045 & +0.004 & [-0.003, +0.011] & 0.2372 \\
test & H1 Hop-PPO - A4 unrestricted & 2045 & +0.010 & [+0.001, +0.018] & 0.0330 \\
ood\_dev & H1 Hop-PPO - M0 oracle & 500 & -0.500 & [-0.542, -0.458] & 0.0000 \\
ood\_dev & H1 Hop-PPO - B1 SFT & 500 & -0.201 & [-0.239, -0.163] & 0.0000 \\
ood\_dev & H1 Hop-PPO - B2a GRPO & 500 & -0.016 & [-0.031, -0.001] & 0.0306 \\
ood\_dev & H1 Hop-PPO - B2b GRPO+path & 500 & +0.020 & [-0.000, +0.041] & 0.0554 \\
ood\_dev & H1 Hop-PPO - B3 soft & 500 & -0.025 & [-0.051, +0.000] & 0.0530 \\
ood\_dev & H1 Hop-PPO - A1 terminal-only & 500 & -0.018 & [-0.036, +0.000] & 0.0532 \\
ood\_dev & H1 Hop-PPO - A2 no-redund. & 500 & -0.049 & [-0.068, -0.029] & 0.0000 \\
ood\_dev & H1 Hop-PPO - A3 no-schema & 500 & +0.006 & [-0.007, +0.019] & 0.3830 \\
ood\_dev & H1 Hop-PPO - A4 unrestricted & 500 & +0.004 & [-0.013, +0.021] & 0.6356 \\
ood\_test & H1 Hop-PPO - M0 oracle & 569 & -0.441 & [-0.480, -0.401] & 0.0000 \\
ood\_test & H1 Hop-PPO - B1 SFT & 569 & -0.141 & [-0.175, -0.107] & 0.0000 \\
ood\_test & H1 Hop-PPO - B2a GRPO & 569 & -0.008 & [-0.022, +0.006] & 0.2672 \\
ood\_test & H1 Hop-PPO - B2b GRPO+path & 569 & +0.025 & [+0.006, +0.045] & 0.0122 \\
ood\_test & H1 Hop-PPO - B3 soft & 569 & -0.006 & [-0.032, +0.019] & 0.6198 \\
ood\_test & H1 Hop-PPO - A1 terminal-only & 569 & -0.009 & [-0.023, +0.004] & 0.1782 \\
ood\_test & H1 Hop-PPO - A2 no-redund. & 569 & -0.028 & [-0.045, -0.012] & 0.0008 \\
ood\_test & H1 Hop-PPO - A3 no-schema & 569 & +0.016 & [+0.005, +0.027] & 0.0062 \\
ood\_test & H1 Hop-PPO - A4 unrestricted & 569 & +0.005 & [-0.009, +0.021] & 0.4834 \\
\bottomrule
\end{tabular}
}
\caption{Paired bootstrap of the EM difference between H1 and each other method on the same questions (10{,}000 resamples): difference, 95\% interval and two-sided $p$.}
\label{tab:paired}
\end{table}

\subsection{Reward and action-space ablations (E2)}
\begin{table*}[t]
\centering\small
\resizebox{\textwidth}{!}{\begin{tabular}{lllllllll}
\toprule
split & method & seeds & EM mean$\pm$std & EM 95\% CI & F1 & strict validity & hop eff & mean hops \\
\midrule
dev & H1 Hop-PPO & 3 & 0.102$\pm$0.012 & [0.091, 0.113] & 0.144$\pm$0.010 & 0.140$\pm$0.044 & 0.078$\pm$0.022 & 2.69$\pm$1.13 \\
dev & A1 terminal-only & 3 & 0.134$\pm$0.060 & [0.122, 0.146] & 0.177$\pm$0.063 & 0.117$\pm$0.035 & 0.072$\pm$0.043 & 4.00$\pm$0.00 \\
dev & A2 no-redund. & 3 & 0.155$\pm$0.041 & [0.142, 0.168] & 0.199$\pm$0.045 & 0.172$\pm$0.025 & 0.102$\pm$0.008 & 3.30$\pm$1.22 \\
dev & A3 no-schema & 3 & 0.106$\pm$0.003 & [0.095, 0.117] & 0.149$\pm$0.006 & 0.131$\pm$0.060 & 0.070$\pm$0.018 & 3.11$\pm$1.05 \\
dev & A4 unrestricted & 3 & 0.102$\pm$0.015 & [0.092, 0.113] & 0.147$\pm$0.018 & 0.097$\pm$0.012 & 0.052$\pm$0.010 & 4.00$\pm$0.00 \\
test & H1 Hop-PPO & 3 & 0.095$\pm$0.013 & [0.088, 0.103] & 0.135$\pm$0.013 & 0.140$\pm$0.044 & 0.069$\pm$0.011 & 2.70$\pm$1.13 \\
test & A1 terminal-only & 3 & 0.119$\pm$0.061 & [0.110, 0.127] & 0.158$\pm$0.069 & 0.121$\pm$0.044 & 0.064$\pm$0.045 & 4.00$\pm$0.00 \\
test & A2 no-redund. & 3 & 0.148$\pm$0.065 & [0.139, 0.156] & 0.194$\pm$0.072 & 0.183$\pm$0.017 & 0.095$\pm$0.031 & 3.30$\pm$1.22 \\
test & A3 no-schema & 3 & 0.091$\pm$0.020 & [0.084, 0.098] & 0.133$\pm$0.023 & 0.134$\pm$0.065 & 0.057$\pm$0.005 & 3.10$\pm$1.05 \\
test & A4 unrestricted & 3 & 0.086$\pm$0.007 & [0.079, 0.092] & 0.129$\pm$0.007 & 0.100$\pm$0.006 & 0.042$\pm$0.004 & 4.00$\pm$0.00 \\
ood\_dev & H1 Hop-PPO & 3 & 0.090$\pm$0.009 & [0.076, 0.105] & 0.145$\pm$0.012 & 0.152$\pm$0.046 & 0.065$\pm$0.015 & 2.70$\pm$1.12 \\
ood\_dev & A1 terminal-only & 3 & 0.108$\pm$0.045 & [0.093, 0.124] & 0.163$\pm$0.049 & 0.126$\pm$0.042 & 0.059$\pm$0.034 & 4.00$\pm$0.00 \\
ood\_dev & A2 no-redund. & 3 & 0.139$\pm$0.039 & [0.121, 0.156] & 0.197$\pm$0.038 & 0.181$\pm$0.020 & 0.089$\pm$0.014 & 3.31$\pm$1.20 \\
ood\_dev & A3 no-schema & 3 & 0.084$\pm$0.005 & [0.070, 0.098] & 0.140$\pm$0.007 & 0.139$\pm$0.061 & 0.054$\pm$0.015 & 3.15$\pm$1.05 \\
ood\_dev & A4 unrestricted & 3 & 0.086$\pm$0.011 & [0.073, 0.100] & 0.145$\pm$0.004 & 0.103$\pm$0.006 & 0.042$\pm$0.003 & 4.00$\pm$0.00 \\
ood\_test & H1 Hop-PPO & 3 & 0.090$\pm$0.018 & [0.076, 0.103] & 0.135$\pm$0.028 & 0.158$\pm$0.059 & 0.074$\pm$0.030 & 2.70$\pm$1.12 \\
ood\_test & A1 terminal-only & 3 & 0.099$\pm$0.045 & [0.085, 0.114] & 0.146$\pm$0.047 & 0.126$\pm$0.040 & 0.057$\pm$0.032 & 4.00$\pm$0.00 \\
ood\_test & A2 no-redund. & 3 & 0.118$\pm$0.039 & [0.103, 0.133] & 0.171$\pm$0.044 & 0.180$\pm$0.015 & 0.080$\pm$0.013 & 3.32$\pm$1.18 \\
ood\_test & A3 no-schema & 3 & 0.074$\pm$0.006 & [0.062, 0.087] & 0.119$\pm$0.007 & 0.140$\pm$0.066 & 0.051$\pm$0.011 & 3.18$\pm$1.01 \\
ood\_test & A4 unrestricted & 3 & 0.084$\pm$0.009 & [0.071, 0.098] & 0.132$\pm$0.012 & 0.109$\pm$0.009 & 0.045$\pm$0.005 & 4.00$\pm$0.00 \\
\bottomrule
\end{tabular}
}
\caption{Ablations of H1 on HybridQA: terminal-only reward (A1), no redundancy penalty (A2), no schema term (A3) and unrestricted action space (A4). Columns as in Table~\ref{tab:main}.}
\label{tab:ablations}
\end{table*}

Removing the hop-level terms (A1, terminal reward only) \emph{raises} EM to \R{EmAOneTerminal} (H1 minus A1: \D{AOneTerminal}) with similar strict validity (\R{ValidAOneTerminal}\% versus \R{ValidHOneHoppo}\%).
The mechanism is visible in the hop counts: with hop rewards and a step cost the policy learns to stop after about two hops (\R{HopsHOneHoppo}), whereas A1 uses all four (\R{HopsAOneTerminal}) and so hands the frozen reader more rows and passages.
H1c (the dense reward matters beyond discrete actions) is therefore not supported for accuracy; the dense reward buys shorter, tidier traces, not more correct answers.
Removing the redundancy penalty (A2, \R{EmATwoNoredund} EM, \D{ATwoNoredund}) has a similar effect for the same reason, while removing the schema term (A3, \D{AThreeNoschema}) or the action-space restriction (A4, \D{AFourUnrestricted}) makes no reliable difference.

\subsection{Learning dynamics and credit assignment (E4)}
\begin{figure*}[t]
\centering
\figfile{F3_learning_curves}{width=0.49\textwidth}\hfill
\figfile{F4_hop_validity}{width=0.49\textwidth}
\caption{Left: development-subset EM versus training trajectories (mean over seeds with interval) for H1 and the trace-level baselines. Right: strict path validity by hop index.}
\label{fig:curves}
\end{figure*}

The learning curves (Figure~\ref{fig:curves}, left) show no faster learning for \hop\ than for the trace-level baselines: all RL methods stay within a few EM points of one another throughout training and well below B1.
Per-hop validity drops sharply after the first hop for every method, so we find no evidence that hop-level credit reduces error propagation in EM terms (H1 inconclusive); the reward-component traces (Appendix, Figure~\ref{fig:reward}) show the hop reward is learned early and the step cost falls once the policy learns to stop.

\subsection{Out-of-domain generalization (E3)}
\begin{table}[t]
\centering\small
\resizebox{\columnwidth}{!}{\begin{tabular}{llllll}
\toprule
pair & method & EM ID & EM OOD & gap abs & gap rel \\
\midrule
\bottomrule
\end{tabular}
}
\caption{EM on in-domain test versus the held-out-cluster OOD split, with absolute and relative gap, per method.}
\label{tab:ood}
\end{table}

\ifcsname GapHOneHoppoOodTest\endcsname
The ID$\rightarrow$OOD gap of \hop\ is \GapHOneHoppoOodTest\ EM points, against \GapBOneSftOodTest\ for B1 and \GapBTwoAGrpoOodTest\ for B2a (Table~\ref{tab:ood}); the OOD split holds out whole table clusters of HybridQA and is not a cross-dataset test.
\else
The OOD evaluation on held-out table clusters (a within-HybridQA split, not a cross-dataset test) is reported in Table~\ref{tab:ood} once every final run has been evaluated; H3 is not assessed in this version.
\fi


